% Archived version DOI: https://doi.org/10.5281/zenodo.22675435
% Living citation annotation; the immutable v0.1.0 PDF and tagged source are preserved.
\documentclass[11pt]{article}
\usepackage[utf8]{inputenc}
\usepackage{amsmath,amssymb,amsthm,mathtools,booktabs}
\usepackage[margin=1in]{geometry}
\usepackage[colorlinks=true,linkcolor=blue,citecolor=blue,urlcolor=blue]{hyperref}
\usepackage{microtype}
\hypersetup{pdftitle={Prefix-width rigidity and dual geometry of multi-marginal earth moving},pdfauthor={Jeffery Kline},pdfsubject={Research preprint, version 0.1.0}}
\pdfinfoomitdate=1
\pdftrailerid{}
\pdfsuppressptexinfo=15
\newtheorem{theorem}{Theorem}[section]
\newtheorem{proposition}[theorem]{Proposition}
\newtheorem{lemma}[theorem]{Lemma}
\newtheorem{corollary}[theorem]{Corollary}
\theoremstyle{definition}\newtheorem{definition}[theorem]{Definition}
\newtheorem{remark}[theorem]{Remark}
\newcommand{\R}{\mathbb R}
\newcommand{\1}{\mathbf 1}
\newcommand{\conv}{\operatorname{conv}}
\newcommand{\aff}{\operatorname{aff}}
\newcommand{\relint}{\operatorname{relint}}
\newcommand{\Phiw}{\Phi_w}
\newcommand{\width}{\operatorname{width}}
\newcommand{\supp}{\operatorname{supp}}
\title{Prefix-width rigidity and dual geometry\\of multi-marginal earth moving}
\author{Jeffery Kline}
\date{September 8, 2026\\\small Version 0.1.0}
\begin{document}
\maketitle
\begin{abstract}
The earth mover's problem asks how cheaply one distribution of mass can be
rearranged into another. Here several histograms on a line are matched at
once, and each unit of matched mass pays the distance between the leftmost and
rightmost bins it occupies. Kline (2019) showed that the minimum cost depends
only on the convex hull of the histograms and adds under Minkowski sums. Both
facts, and explicit optimal primal and dual solutions, follow from one formula:
at each of the $n-1$ cuts between adjacent bins, take the largest cumulative
mass to the left of the cut minus the smallest, and add these gaps over all
cuts. The main result is a converse. Consider any rule that assigns a cost to
each tuple of bins, for any number of histograms, and is symmetric,
nonnegative, zero for a single histogram, and Monge, the submodularity
condition under which greedy matching is optimal. Its minimum cost is hull
invariant and Minkowski additive for every common mass if and only if the rule
is a nonnegative weighted count of the cuts lying between the leftmost and
rightmost bins. Charging one per cut recovers the range cost exactly. Two
further results describe the dual geometry. After quotienting out the
lineality space, the shifts of the dual potentials that change nothing, the
dual polyhedron for $d$ histograms on $n$ bins has exactly $[d(d-1)]^{n-1}$
vertices, all integral, one for each choice of a lowest and a highest histogram
at every cut. When those choices form a star with at least two leaves,
realizing them needs, and is achieved in, affine dimension two.
\end{abstract}

\paragraph{Scope and priority.}
The classification concerns the fixed-order, all-mass, all-arity transport
class stated below. The source comparison recorded with this preprint did not
identify either headline theorem in the inspected literature; it does not
establish historical originality. The star-dimension construction is an
explicit example, with no claim of originality.

\section{Introduction}
The earth mover's problem asks how cheaply one distribution of mass can be
rearranged into another. This paper concerns a version with several histograms
on a line. Each histogram places the same total mass in $n$ ordered bins. A
coupling assigns mass to tuples containing one bin from each histogram, so that
every histogram's bins receive their prescribed masses. A tuple costs the
distance between its leftmost and rightmost bins, and the value of the problem
is the least total cost of a coupling. Kline~\cite{kline} proved that this
value is homogeneous, invariant under translating every histogram by the same
vector, monotone under enlargement of the convex hull of the histograms, and
additive under Minkowski sums.

The value has a simple description in cumulative masses. At each of the $n-1$
cuts between adjacent bins, record how much mass each histogram places to the
left. Subtract the smallest of these amounts from the largest and add the
results over all cuts. This sum is exactly the minimum cost
(Theorem~\ref{thm:width}). In geometric terms, each summand is the width of the
histograms' convex hull in a prefix direction; in cumulative coordinates, the
value is the sum of the side lengths of the smallest box containing the
histograms.

The four properties follow at once. Enlarging the convex hull can only
increase its widths. Translating every histogram by the same vector leaves the
widths unchanged whenever the translated histograms remain admissible. For two
histogram families $X$ and $Y$, the Minkowski sum $X+Y$ consists of every
pairwise sum $x+y$, with repetition; its common mass is the sum of the two
common masses, and each of its prefix widths is the sum of the corresponding
widths of $X$ and $Y$, so the values add. This is the Minkowski additivity used
throughout.

The main result asks which other costs behave this way, and the answer is that,
up to reweighting the cuts, none do. Fix the bin order and a cost for every
number of histograms. Require each cost to be finite, nonnegative, unchanged by
permuting the histograms, and Monge, the submodularity condition under which
greedy matching is optimal, with zero cost for a single histogram. Suppose the
value depends only on the convex hull of the histograms and is Minkowski
additive, for every common mass. Theorem~\ref{thm:rigidity} shows that the cost
must then be a nonnegative weighted count of the cuts lying between the
leftmost and rightmost occupied bins. Only $n-1$ weights remain, one per cut,
and each is the two-histogram cost between the adjacent bins at that cut.
Charging one per cut fixes the whole cost family and recovers range-cost EMD.
Section~4 states the assumptions precisely and proves both directions.

This characterization sits inside a broad geometric theory of additive
functionals. Nonnegative sums and averages of directional widths, mean width
among them, are all Minkowski additive. In cumulative coordinates, range-cost
EMD is the standard $\ell_1$ diversity, and Bryant and
Tupper~\cite[Section 2.1 and Theorem 5]{bt} describe the linear semidiversities
as support integrals against positive measures with zero first moment. What
the transport hypotheses add is rigidity: they force the particular prefix
directions and leave only their weights free. The uniqueness claim is confined
to this transport class; the mean-width example in Section~4 marks the
boundary.

The width formula also yields explicit optimal solutions of the transport
problem and of its dual. We use the classical optimality of common-quantile
couplings for Monge costs and include an uncrossing proof. Queyranne, Spieksma
and Tardella~\cite[Theorem 3.1(2)--(3)]{qst} give a general greedy primal/dual framework. Dual
solutions serve as gradients in the machine-learning applications of Mehta,
Kline, Lokhande, Fung and Singh~\cite[Theorem 4.3 and Appendix 7.1]{mehta}; the prefix formula makes them
explicit and separates derivatives at strict extrema from subgradients at ties.

Theorem~\ref{thm:census} then describes every vertex of the dual for $d$
histograms. A polyhedron's \emph{lineality space} consists of the directions in
which one can move arbitrarily far both forward and backward without leaving
it. In our dual, adding a constant $\alpha_j$ to every entry of the potential
$y_j$ changes neither feasibility nor the objective when $\sum_j\alpha_j=0$,
because each constraint uses one entry from every potential and the histograms
have equal mass. These shifts are the lineality space, and working
\emph{modulo lineality} means identifying solutions that differ only by such a
shift. Before this quotient the dual has no vertices at all when $d\ge2$, since
every feasible point lies on a feasible line. After it, the vertices are
counted exactly: one for each choice of a lowest and a highest histogram at
every cut. Section~3 gives the quotient and the count.

We also ask how many dimensions are needed to realize prescribed pairs of
histograms as the unique lowest and highest at successive cuts. This is a
strict-antipodality problem in convex geometry. For a star of prescribed pairs
with at least two leaves, two dimensions suffice and one is impossible.
Section~\ref{sec:branch} turns to a different problem, a star as the underlying
transport space rather than a line; its three-marginal formula follows from
Makur and Singh~\cite{makur} and shows why branching can prevent all cut bounds
from being attained at once. The appendix records qualifications to the 2019
algorithm and to one perturbation argument.

\section{Histograms, cuts and exact certificates}
Fix $n\ge2$ bins indexed by $0,\ldots,n-1$. A histogram is a vector in
$\R_+^n$, where $\R_+$ includes zero. Let $X=(x_1,\ldots,x_d)$ be a nonempty
finite multiset of histograms with common mass $m$. A coupling is an array
$f\colon\{0,\ldots,n-1\}^d\to\R_+$ with marginals $x_j$. For a cost $c_d$,
write
\[
 \Phi_c(X)=\min_f\sum_I c_d(I)f(I).
\]
The feasible set is nonempty and compact. For the range cost
$c_d(I)=\max_j I_j-\min_j I_j$, write $\Phi(X)$. For $d=1$ this cost is zero.

Put $H_m=\{x:\1^Tx=m\}$ and define the prefix map
\[
 F_x(k)=\sum_{i=0}^k x(i),\qquad
 T(x)=(F_x(0),\ldots,F_x(n-2)).
\]
Restricted to $H_m$, $T$ is an affine bijection onto $\R^{n-1}$: successive
differences recover $x$, including $x(n-1)=m-F_x(n-2)$. Let
$u_k=(\underbrace{1,\ldots,1}_{k+1},0,\ldots,0)$. For nonempty compact $K$,
$h_K(u)=\max_{x\in K}\langle u,x\rangle$ and
$\width(K,u)=h_K(u)+h_K(-u)$.

\begin{theorem}[Prefix widths]\label{thm:width}
For every admissible $X$,
\begin{equation}\label{eq:width}
 \Phi(X)=\sum_{k=0}^{n-2}\left(\max_j F_{x_j}(k)-\min_j F_{x_j}(k)\right)
        =\sum_{k=0}^{n-2}\width(\conv X,u_k).
\end{equation}
The common-quantile coupling attains the minimum, including at zero bins and ties.
\end{theorem}
\begin{proof}
For every tuple $I$,
\[
 \max I-\min I=\sum_{k=0}^{n-2}\1\{\min I\le k<\max I\}.
\]
In a coupling of mass $m$, let $A_{j,k}=\{I_j\le k\}$; this event has measure
$F_{x_j}(k)$. The measure of the event in the $k$th summand is
$\mu(\bigcup_j A_{j,k})-\mu(\bigcap_j A_{j,k})$, at least
$\max_j F_{x_j}(k)-\min_j F_{x_j}(k)$.
For $t\in[0,m]$, let $I_j(t)$ be the quantile bin of $x_j$. Then, up to null
endpoints, $A_{j,k}=[0,F_{x_j}(k)]$. All the cut bounds are attained
simultaneously. Integration proves the result. If $m=0$, all terms vanish.
\end{proof}

\begin{proposition}[An explicit dual]\label{prop:dual}
For each cut $k$, choose $a_k$ minimizing and $b_k$ maximizing $F_{x_j}(k)$.
Set
\[
 q_k=e_{a_k}-e_{b_k},\qquad y_j(i)=\sum_{k<i}q_k(j).
\]
Here $q_k=0$ is permitted when the extrema tie. Then $y$ is dual feasible and
$\sum_j\langle x_j,y_j\rangle=\Phi(X)$. Its entries are integers bounded in
absolute value by $n-1$, every row sums to zero, and at most $2(n-1)$ columns
are nonzero.
\end{proposition}
\begin{proof}
For every tuple $I$,
\[
 \sum_j y_j(I_j)=\sum_k\sum_{j:I_j>k}q_k(j)\le \max I-\min I.
\]
Cuts outside $[\min I,\max I)$ contribute zero; each remaining cut contributes
at most one. Summation by parts gives
\[
 \sum_j\langle x_j,y_j\rangle
 =\sum_k\sum_j q_k(j)(m-F_{x_j}(k))
 =\sum_k\bigl(F_{x_{b_k}}(k)-F_{x_{a_k}}(k)\bigr).
\]
The remaining assertions follow from the root increments.
\end{proof}

The quantile coupling has at most $d(n-1)+1$ positive cells. The value and dual
can be obtained in $O(nd)$ arithmetic operations by accumulating prefixes and
extrema. Producing the primal coupling additionally requires merging the
breakpoints and reporting its cells; the value-only complexity is not a claim
about that separate output task.

\begin{corollary}[Chambers and sensitivity]\label{cor:gradient}
On the relative interior of a fixed-mass histogram domain, if each prefix has
unique minimizing and maximizing labels, the functional is locally linear and
the dual in Proposition~\ref{prop:dual} represents its derivative. At ties,
any certificate from that proposition is a subgradient. A chosen certificate
uses at most $2(n-1)$ marginal columns.
\end{corollary}
\begin{proof}
Strict comparisons persist in a neighborhood, where the width formula uses
the same linear summands. At ties, dual feasibility for every admissible $X'$ and
equality at $X$ give
$\Phi(X')\ge\langle X',y\rangle=\Phi(X)+\langle X'-X,y\rangle$.
The column bound was already proved. Derivatives are understood on the
fixed-mass tangent space, so columnwise additive constants are immaterial.
\end{proof}

This is an explicit certificate interpretation of the gradient use in
\cite[Theorem 4.3 and Appendix 7.1]{mehta}, not a claim that differentiating transport optima is new. It also
makes the tie qualification visible: a subgradient need not be an ordinary
derivative. A point can lie outside the current convex hull while making no
incremental contribution, as the interval formula below shows.

\subsection{Geometric consequences and limitations}
\begin{corollary}\label{cor:geometry}
The range functional depends only on $\conv X$, is positively homogeneous,
is invariant under admissible common translations, and is monotone under
convex-hull inclusion. For multiset Minkowski addition,
\[
 \Phi(X+Y)=\Phi(X)+\Phi(Y),
\]
where $X+Y=(x_j+y_l)_{j,l}$ retains multiplicities and its common mass is the
sum of the two common masses. Moreover,
$\Phi(\{x,y\})=\|T(x-y)\|_1$ on a common mass hyperplane.
\end{corollary}
\begin{proof}
All statements follow from~\eqref{eq:width}. For additivity use
$h_{K+L}=h_K+h_L$. For the two-point formula, a width is the absolute difference
of the two corresponding prefixes. Injectivity of $T$ on $H_0$ gives separation.
\end{proof}

If $L\subset H_m$ is compact and convex, $K$ is nonempty and compact,
$K\subset\relint L$, and $\dim L>0$, then the
width sum is strictly smaller on $K$. Every linear functional nonconstant on
$\aff L$ takes its extrema on the relative boundary, so its extrema on $K$ are
strictly interior to its range on $L$. At least one prefix functional is
nonconstant because the prefixes separate points of $H_m$. Positive dimension
and nonempty sets exclude the singleton and empty-set degeneracies.

\begin{proposition}[Exact witnesses and increments]\label{prop:witness}
For every finite $X$ there is a submultiset of at most $2(n-1)$ points with the
same value. This bound is worst-case sharp. With
$L_k=\min_X F_x(k)$ and $U_k=\max_X F_x(k)$, a further histogram $z\in H_m$ has
\[
 \Phi(X\cup\{z\})-\Phi(X)
 =\sum_{k=0}^{n-2}\operatorname{dist}\bigl(F_z(k),[L_k,U_k]\bigr).
\]
\end{proposition}
\begin{proof}
Choose one point attaining each signed prefix extremum. For sharpness, place a
small translated crosspolytope with vertices $\pm e_k$ in the interior of the
prefix image of the histogram simplex. Each signed coordinate extremum is
unique, so all $2(n-1)$ vertices are needed. The increment formula follows by
expanding the interval length after adjoining one value.
\end{proof}

In particular, zero incremental contribution need not imply membership in
$\conv X$. The functional retains the coordinate bounding box of $T(X)$ and
forgets correlations among its coordinates. This limits deductions about shape,
lattice points or arithmetic solvability that rely on the value alone.
For compact convex sets the same width sum is defined directly by supports;
choosing signed support points gives the same finite witness bound. If $C$ lies
in the histogram simplex, the polytope
\[
 \{x\in H_m\cap\R_+^n:L_k\le F_x(k)\le U_k\ \text{for every }k\}
\]
contains $C$ and has exactly its prefix extrema. This gives the enclosing-polytope
value directly, without a limiting argument.

\section{The complete dual vertex set}
For $n,d\ge2$, define the full dual polyhedron
\[
 D_{n,d}=\{z\in\R^{nd}:\textstyle\sum_j z_j(I_j)\le\max I-\min I\text{ for all }I\}.
\]
Its lineality is
\[
 L=\{z:z_j(i)=\alpha_j\text{ for all }i,\ \textstyle\sum_j\alpha_j=0\}.
\]
Indeed a lineality direction has zero sum on every tuple; varying one index
makes each column constant. Thus $D_{n,d}/L$ has dimension $nd-d+1$.
Without this quotient the polyhedron has no vertices.

\begin{theorem}[Full dual census]\label{thm:census}
Let $P\subset D_{n,d}$ be the face where all diagonal constraints are tight.
Modulo $L$, it is linearly isomorphic to
\[
 R_d^{\,n-1},\qquad R_d=\conv\{e_a-e_b:a\ne b\}.
\]
Every vertex of $D_{n,d}/L$ lies in $P/L$. Consequently the full quotient dual
has exactly $[d(d-1)]^{n-1}$ vertices.
\end{theorem}
\begin{proof}
The simultaneous diagonal equality set is a face: the sum of its nonnegative
diagonal slacks vanishes exactly when every slack vanishes. On this face there
is a unique representative with $z_j(0)=0$ for all $j$. Write
$q_k(j)=z_j(k+1)-z_j(k)$. Diagonal equalities give $\sum_jq_k(j)=0$.
Tuples using only $k,k+1$ give
\[
 \sum_{j\in S}q_k(j)\le1\quad(\varnothing\ne S\subsetneq\{1,\ldots,d\}).
\]
Equivalently, $\sum_j(q_k(j))_+\le1$. This describes $R_d$: decompose the
positive and negative parts into a nonnegative flow of total mass at most one,
and use $0\in R_d$. Conversely, arbitrary $q_k\in R_d$ give a feasible dual
by the cut argument in Proposition~\ref{prop:dual}. The cuts are independent,
so $P/L\cong R_d^{n-1}$. Every root is exposed by a functional with unique
largest coordinate at its positive endpoint and unique smallest coordinate at
its negative endpoint. This proves the vertex count for the face.

To prove completeness, take a vertex $v$ of the full quotient. Its active tuple
normals span $L^\perp$. Let $x$ be their sum. Each normal is nonnegative and
has column sums one. Every coordinate of $x$ is positive: otherwise the active
normals would all lie in a proper coordinate hyperplane of $L^\perp$, contrary
to spanning. Thus $x$ is an admissible family of strictly positive histograms.
The objective $x$ uniquely exposes $v$ modulo $L$, since attaining its upper
bound requires equality in every active inequality. But
Proposition~\ref{prop:dual} supplies an optimal point in $P$ for this same
objective. Uniqueness puts $v$ in $P/L$. Finally a vertex of a face is a vertex
of the full polyhedron, proving the assertion.
\end{proof}

\begin{remark}
The bounded product is not the whole quotient. The recession cone is
\[
 \operatorname{rec}D_{n,d}
 =\{r:\textstyle\sum_j\max_i r_j(i)\le0\}
 =-\R_+^{nd}+L.
\]
For the second equality, choose constants $\alpha_j\ge\max_i r_j(i)$ whose sum
is zero. The vertex--recession decomposition then gives
$D_{n,d}=P+\operatorname{rec}D_{n,d}$.
\end{remark}

\section{Rigidity of normalized additive Monge families}
\subsection{The all-arity problem}
Fix the bin order once and for all. For every $d\ge1$, let
$c_d:\{0,\ldots,n-1\}^d\to\R_+$ be finite and symmetric. A cost is Monge if
\[
 c_d(I\wedge J)+c_d(I\vee J)\le c_d(I)+c_d(J),
\]
with coordinatewise meet and join. Normalization means $c_1(i)=0$ for every bin.
Hull invariance of the associated family of transport functionals means that
$\conv X=\conv Y$ implies $\Phi_c(X)=\Phi_c(Y)$, for any admissible multisets
in the same mass hyperplane, including different cardinalities. Additivity means
\[
 \Phi_{c_{de}}(X+Y)=\Phi_{c_d}(X)+\Phi_{c_e}(Y).
\]
The masses of $X$ and $Y$ may differ. No comparison of differently shaped cost
tensors is made without this functional interpretation.

\begin{lemma}[Common-quantile optimality]\label{lem:monge}
A finite Monge cost has an optimal common-quantile coupling for every family of
equal-mass nonnegative histograms.
\end{lemma}
\begin{proof}
Among optimal couplings, maximize the secondary linear objective
$\sum_I f(I)(\sum_j I_j)^2$. Both extrema exist by compactness. If two
incomparable tuples $I,J$ have positive mass, replace an equal small amount of
mass at them by mass at $I\wedge J,I\vee J$. Each marginal is preserved.
The Monge inequality prevents the primary objective from increasing; optimality
therefore makes its change zero. Put $a=\sum I_j$, $b=\sum J_j$ and
$t=\sum\min(I_j,J_j)$. The secondary change per unit mass is
\[
 t^2+(a+b-t)^2-a^2-b^2=2(a-t)(b-t)>0,
\]
contradiction. Hence the support is a chain. Order its atoms and assign
consecutive intervals of their masses on $[0,m]$. Every coordinate is
nondecreasing and has the prescribed marginal, so this is the common-quantile
coupling, up to null endpoints.
\end{proof}

\begin{theorem}[Normalized classification]\label{thm:rigidity}
A symmetric, finite, nonnegative Monge family $(c_d)_{d\ge1}$ with $c_1=0$
has a hull-invariant, Minkowski-additive transport functional if and only if
there are weights $w_0,\ldots,w_{n-2}\ge0$, independent of $d$, such that
\begin{equation}\label{eq:weightedcost}
 c_d(I)=\sum_{k=0}^{n-2}w_k\1\{\min I\le k<\max I\}.
\end{equation}
The weights are uniquely determined by $w_k=c_2(k,k+1)$, and
\[
 \Phi_c(X)=\sum_{k=0}^{n-2}w_k\width(\conv X,u_k).
\]
\end{theorem}
\begin{proof}
For Dirac histograms $e_{I_1},\ldots,e_{I_d}$ the only feasible coupling is the
unit atom at $I$, so its optimum is $c_d(I)$. The convex hull of these histograms
depends only on the set of bin labels occurring in $I$. Hull invariance therefore
implies that repeating labels does not change this cost. In particular every
constant tuple has cost zero, and every mixed tuple using just $k,k+1$ has cost
$w_k=c_2(k,k+1)$, independent of arity and multiplicities.

Take arbitrary $X$ of mass $m$ and choose $M>m$. Add the common histogram
$h=M\1$. Normalization and Minkowski additivity give
\begin{equation}\label{eq:padding}
 \Phi_c(X+h)=\Phi_c(X)+\Phi_{c_1}(\{h\})=\Phi_c(X).
\end{equation}
By Lemma~\ref{lem:monge}, the translated family has an optimal common-quantile
coupling. Its cut-$k$ breakpoints are
\[
 (k+1)M+F_{x_j}(k).
\]
The $n-1$ breakpoint bands lie in the disjoint intervals
$[(k+1)M,(k+1)M+m]$. Outside the mixed portions of these bands the index tuple
is constant, hence costs zero. Inside band $k$, the only possible bins are
$k,k+1$, and the tuple is mixed exactly between its least and greatest
breakpoint. The cost there is $w_k$. Thus
\[
 \Phi_c(X+h)=\sum_k w_k\bigl(\max_jF_{x_j}(k)-\min_jF_{x_j}(k)\bigr).
\]
Together with~\eqref{eq:padding} this gives the asserted functional formula.
Applying it to the Dirac family recovers~\eqref{eq:weightedcost}. Nonnegativity
of $c_2$ gives $w_k\ge0$.

Conversely, a cut indicator is submodular: under $S=\{j:I_j>k\}$ it becomes
$g(S)=1$ for nonempty proper $S$ and zero otherwise. Its cardinality increments
are nonincreasing, so it is submodular; meet and join map to intersection and
union. Nonnegative weighted sums are consequently Monge. The weighted version
of Theorem~\ref{thm:width} proves the functional formula and hence hull
invariance and Minkowski additivity. Symmetry and normalization are immediate.
\end{proof}

\paragraph{Unit adjacent costs.}
Within the class of Theorem~\ref{thm:rigidity}, the additional conditions
$c_2(k,k+1)=1$ for every $k$ uniquely determine $c_d(I)=\max I-\min I$ for
all $d$. Indeed the theorem identifies these adjacent costs with the weights
in~\eqref{eq:weightedcost}, so every weight is one. The range family satisfies
all the assumptions by Theorem~\ref{thm:width}; this gives existence as well
as uniqueness.

This conclusion does not extend to arbitrary additive geometric functionals.
For three bins, the Dirac histograms have prefix coordinates $(1,1),(0,1),(0,0)$.
Mean width normalized to induce Euclidean distance on pairs~\cite[Section 2.1]{bt}
gives adjacent pair values $1,1$ but endpoint value $\sqrt2$, rather than $2$.
A transport representation would inherit these Dirac costs, and its Monge
inequality for $I=(0,2)$, $J=(1,1)$ would require $2\le\sqrt2$. Thus this
additive alternative cannot have a transport representation satisfying the
ordered Monge requirement.

\begin{remark}[What the hypotheses do]
The proof needs hull invariance only to remove repeated Dirac labels, and
additivity only for common uniform padding. Thus those two weaker invariances,
together with the other stated assumptions, already force the conclusion.
The bin order is fixed; the statement is not a classification over arbitrary
underlying metric spaces. If only probability histograms are used, a separate
normalization of the sum operation is needed; that different formulation is not
silently assumed here.
\end{remark}

\subsection{Controls on the assumptions}
The following examples show that additivity and the fixed-order Monge condition
are substantive restrictions.

\begin{proposition}[Hull invariance without additivity]\label{prop:squared}
The normalized family $c_d(I)=(\max I-\min I)^2$ is symmetric, nonnegative and
Monge. Its transport functional is hull invariant, but not Minkowski additive.
Writing $L_k=\min_j F_{x_j}(k)$ and $U_k=\max_j F_{x_j}(k)$, its value is
\[
 \Phi_c(X)=\sum_k(U_k-L_k)+2\sum_{k<l}(U_k-L_l)_+.
\]
\end{proposition}
\begin{proof}
To check Monge, write the index intervals of $I,J$ as $[a,A],[b,B]$, with
$a\le b$. If $A\le B$, meet and join have widths at most $A-a$ and $B-b$,
respectively. If $B\le A$, their widths are at most $B-a$ and $A-b$.
These latter lengths have the same sum as $A-a,B-b$, and their sum of squares
is smaller by $2(b-a)(A-B)$. This proves the required inequality in both cases.

In a quantile coupling, crossing cut $k$ corresponds to the interval
$[L_k,U_k]$. Since prefixes increase with $k$, crossing both $k<l$ corresponds
to an interval of length $(U_k-L_l)_+$. Expand the square of the number of
crossed cuts and apply Lemma~\ref{lem:monge}. Every term depends only on the
prefix extrema, proving hull invariance. For $n=3$ and $X=\{e_0,e_2\}$, the
value is $4$, whereas adding $h=2(1,1,1)$ to every histogram gives value $2$.
The singleton $\{h\}$ has value zero, so additivity fails.
\end{proof}

For the Monge restriction, reorder the three bins as $(0,2,1)$ and let the
cost be their range in that reordered path. This normalized family is hull
invariant and additive by the reordered prefix formula. It is not Monge in
the original order: for $I=(0,2)$ and $J=(1,1)$ the meet/join costs sum to
$3$, while the original costs sum to $1$. Thus removing Monge admits even
other path orders, not just different weights on the fixed path.

\subsection{Why normalization matters; a related conjecture}
The unqualified endpoint-only converse is false. The family
$c_d(I)=d^{-1}\sum_j I_j$ is symmetric, nonnegative and modular, hence Monge.
Its transport value is the average unary marginal moment and is additive under
multiset Minkowski addition. But $(0,0,1)$ and $(0,1,1)$ have the same endpoints
and costs $1/3$ and $2/3$. This family is not normalized and is sensitive to
multiplicities; it does not contradict Theorem~\ref{thm:rigidity}.

More generally let $a$ be nondecreasing and $b$ nonincreasing, and set
$c_d(I)=a(\max I)+b(\min I)$. With $L_k=\min_jF_{x_j}(k)$ and
$U_k=\max_jF_{x_j}(k)$, the common quantile simultaneously minimizes the terms
in the following expansion:
\[
 \Phi_c(X)=m[a(n-1)+b(n-1)]-\sum_k\Delta a_k L_k-\sum_k\Delta b_k U_k.
\]
Indeed $\{\max I>k\}$ complements the intersection of prefix events, while
$\{\min I>k\}$ complements their union; the signs of the increments give the
correct respective bounds. The formula proves additivity. In particular
\[
 \Phi_{\max}(X)=m(n-1)-\sum_k L_k,
\]
settling the maximum-cost additivity conjecture in~\cite[Section 5]{kline}.
If every diagonal cost is zero, $a+b=0$, reducing this family to weighted ranges.

\section{Endpoint graphs and minimum affine dimension}
Let $G$ be a finite loopless graph without isolated vertices, with a nonempty
ordered edge list of length $m$. Choose an orientation for each edge and put
$n=m+1$. A strict endpoint realization consists of positive probability
histograms $x_j\in\R^n$ such that the $k$th prefix has its unique minimum
and maximum at the prescribed endpoints of edge $k$. Extra pairs may be
strictly antipodal in other directions; this is not an induced-graph condition.
Let $\rho(G)$ be the least affine dimension of such a realization.

\begin{proposition}[Unrestricted realization]\label{prop:realize}
Every such ordered endpoint pattern has a realization. If $(a_k,b_k)$ are its
minimum and maximum labels, use
\[
 F_{x_j}(k)=\frac{k+1}{n}+\epsilon\delta_{jk},\qquad
 \delta_{jk}=\begin{cases}-1&j=a_k,\\1&j=b_k,\\0&\text{otherwise},\end{cases}
 \quad 0<\epsilon<\frac1{2n}.
\]
The resulting histograms are positive, every label is an exposed vertex, and
the normalized optimal dual is unique with increments $e_{a_k}-e_{b_k}$.
\end{proposition}
\begin{proof}
Successive prefix differences are positive, including the two boundary bins.
The extrema are strict. Every label occurs on an edge and is therefore uniquely
exposed by a signed prefix direction. Consecutive prefix ranges are disjoint,
so the quantile coupling gives positive mass to every diagonal tuple. Any
optimal dual consequently has all diagonals tight. Within the product face in
Theorem~\ref{thm:census}, the strict extrema uniquely expose the stated roots.
\end{proof}

\begin{definition}
Two points $p_a,p_b$ in a finite set are strictly antipodal if a linear functional
is uniquely minimized at one and uniquely maximized at the other over that set.
\end{definition}
\begin{theorem}[Dimension reduction]\label{thm:antipodal}
The number $\rho(G)$ equals the least $r$ for which points $p_j\in\R^r$ can
make every requested edge a strictly antipodal pair. In particular edge order,
orientation and repetitions do not affect $\rho(G)$.
\end{theorem}
\begin{proof}
Restrict the prefix functionals of a histogram realization to its affine hull.
This gives a strictly antipodal realization in the same dimension. Conversely,
choose a witnessing functional $\ell_k$ for each edge of a point realization
and set
\[
 F_{x_j}(k)=\frac{k+1}{m+1}+\epsilon\ell_k(p_j).
\]
There are finitely many values, so sufficiently small positive $\epsilon$ makes
all successive differences and boundary bins positive. The requested extrema
are unchanged. The histogram affine dimension is at most $r$; no label
collapses or ceases to be exposed because every label is incident to a strict
edge. Taking minima in both directions proves equality.
\end{proof}

\begin{theorem}[Stars]\label{thm:stars}
For the star $K_{1,s}$, $\rho(K_{1,1})=1$, and $\rho(K_{1,s})=2$ for every
$s\ge2$.
\end{theorem}
\begin{proof}
A segment realizes the single edge. For $s\ge2$, choose distinct rational
$t_1,\ldots,t_s\in[-1/4,1/4]$, an apex $p_0=(0,1)$, and leaves
$p_j=(t_j,t_j^2)$. The functional $\ell_j(x,y)=y-2t_jx$ satisfies
\[
 \ell_j(p_i)=(t_i-t_j)^2-t_j^2,\qquad \ell_j(p_0)=1.
\]
It is uniquely minimized among the leaves at $p_j$, and every leaf value is
at most $3/16<1$. Hence it is uniquely maximized at the apex. This gives all
star edges as strict antipodal pairs in the plane, and
Theorem~\ref{thm:antipodal} gives the histogram realization. In one dimension
only the two endpoints of a finite set can be uniquely exposed; since all
$s+1\ge3$ labels must be exposed, dimension one is impossible.
\end{proof}

The construction is explicit in histogram coordinates: with $n=s+1$, take
$\epsilon=1/(4n)$ in the proof of Theorem~\ref{thm:antipodal}. The values
$\ell_j(p_i)$ lie in $[-1/16,1]$, ensuring every required prefix difference is
positive. Thus the theorem produces rational histograms, not just an abstract
existence argument.

For contrast, the complete graph requires a pairwise strictly antipodal set.
If $d$ such points span dimension $r$, the copies $(K+p_j)/2$ inside their
convex hull $K$ have disjoint interiors: a functional witnessing an antipodal
pair separates the corresponding two copies. Comparing volumes gives $d\le2^r$.
The construction of Barvinok, Lee and Novik~\cite[Theorem 4.1]{bln} gives at
least $3^{\lfloor r/2\rfloor-1}-1$ pairwise strictly antipodal points. Consequently
\[
 \lceil\log_2d\rceil\le\rho(K_d)
 \le2\lceil\log_3(d+1)\rceil+2.
\]
These are transferred existing bounds, not new extremal estimates. They separate
the constant-dimensional star family from complete graphs. The exact star
example is included to make the dimension parameter concrete; its originality
is not a premise of the paper.

\section{The branching-star connection to coupling theory}\label{sec:branch}
Now let the underlying transportation space be a unit three-leaf star, with
three probability marginals $p_1,p_2,p_3$ supported on its leaves. The cost is
the length of the minimal subtree spanning a sampled tuple. Define
\[
 Q=\sum_v\min_i p_i(v),\quad B=\sum_v\operatorname{med}_i p_i(v),\quad
 H=\sum_v\max_i p_i(v),\quad L=H-Q.
\]
Here $L$ is the sum of the edgewise cut lower bounds. The exact formula is
\begin{equation}\label{eq:star}
 \Phi_{\mathrm{star}}=L+(B-1)_+.
\end{equation}
This is a known corollary, not a new theorem claimed here. To make the
translation explicit, let $D$ be the number of distinct leaves in a tuple.
Pointwise the spanning cost is $D-\1\{\text{all three equal}\}$.
Makur and Singh~\cite[Proposition 1 and Sections III-C1, III-D]{makur} show
that the minimum expected $D$ is $H+(B-1)_+$ and construct an attaining
coupling with diagonal mass $Q$. Every coupling has diagonal mass at most $Q$;
the same coupling attains both bounds, giving~\eqref{eq:star}. Subtracting
unrelated extrema would not suffice.

This connects the convex-width interpretation of~\cite{kline} to the later
coupling result without asserting a historical citation link. The accessible
Makur--Singh manuscript, arXiv version 2, does not cite Kline or the EMD paper.
The local star proof uses the same common-diagonal, pair-agreement and residual
mass mechanism, reorganized as a direct construction. Its independent derivation
does not establish a different method or a novel formula.

For the cyclic half-mass marginals $(0,1/2,1/2)$, $(1/2,0,1/2)$ and
$(1/2,1/2,0)$, $L=3/2$, $B=3/2$ and the optimum is $2$. Thus a common
coupling need not attain all tree-edge bounds. The path proof works because all
prefix events can be nested simultaneously; the obstruction on a branching
space identifies exactly where that proof mechanism can fail.

\section{Evidence, scope and remaining review}
The width and certificate identities were checked on 80 rational cases, with
144 corresponding raw transport LP instances; these instances were repeated
after implementing the reusable solver. Complete rational enumeration of the
raw full-dual inequalities gave respectively $6,4,12$ vertices for
$(n,d)=(2,3),(3,2),(2,4)$. These checks supplement the all-parameter proofs;
they do not prove originality or replace mathematical examination.

The draft-specific checks used 24 rational weighted-cost instances and 122
raw LP solves, testing cross-mass addition, convex-hull augmentation, uniform
padding and the squared-range control. The largest absolute LP discrepancy was
$1.17\cdot10^{-15}$. Exact arithmetic checked 11,467 Monge inequalities for
selected squared-range tensors. The rational star construction was checked for
$2,3,4,7,11,20,32$ leaves, including positivity, strict endpoints and affine
rank two; two further small raw LP solves checked its transport values.
These are bounded checks of examples, distinct from the symbolic proofs.

The normalized classification uses its explicit all-mass, all-arity
assumptions. A separate AI review checked the padding and Dirac reductions,
the full-dual completeness argument, zeros, ties and the star-dimension
construction. The source comparison includes the final conference-formatted
ICLR paper and its technical appendix. These checks are documented in the
accompanying audit and reproduction records; they are not expert peer review
or a certificate of novelty.

AI assisted the development, proof examination, code, source comparison and
exposition of this work under the author's direction. The author is responsible
for the claims and for correcting or withdrawing them when necessary.

Further graph-family classifications may be worthwhile, but are not prerequisites
for the statements proved here. Before claiming a new graph-theoretic result,
one should compare the selected invariant with the established strict-antipodality
literature. Arithmetic applications require additional information beyond the
prefix bounding box and are outside this draft's claims.

\clearpage
\appendix
\section{Qualifications to the 2019 paper}
These observations concern particular statements of~\cite{kline}, not the
validity of the geometric functional, which is proved directly above. They are
kept separate from the proposed classification and dual-geometry results.

\subsection{A zero-bin dual-feasibility failure}
In Algorithm 1 of~\cite[p.136]{kline}, take two identical histograms
$x_0=x_1=(1,0)$ and the stated smallest-label tie rule. The first iteration
removes mass one at $(0,0)$, advances the first index to $(1,0)$ and assigns
$y_0(1)=1$. At zero residual mass the next iteration again advances the first
index, now out of range, and terminates. The resulting dual is
\[
 y_0=(0,1),\qquad y_1=(0,0).
\]
At tuple $(1,1)$ its sum is one, violating the zero-cost diagonal constraint.
Both reported objective values are nevertheless zero because the offending
coordinates have zero marginal mass. The unchanged archived GEM implementation
reproduces this example; its unfinished-suffix fill does not repair it. This
is not a claim about an uninspected current remote revision.

Proposition~\ref{prop:dual} gives a direct replacement valid on the closed
nonnegative orthant. For this example, choosing the same minimizing and maximizing
label yields the identically zero dual. The local implementation computes the
prefix certificate directly and optionally returns the quantile coupling. Any
historical correction to the printed algorithm should be reviewed explicitly.

\subsection{Relative interior and transverse translation}
Theorem 3.1(4) of~\cite[p.136]{kline} requires containment under all sufficiently
small translations in $H_0$. This fails when the containing hull is
lower-dimensional. Take
\[
 Z=\{(1/2,0,1/2),(0,1/2,1/2)\},\quad X=\{(1/4,1/4,1/2)\}.
\]
Then $X\subset\relint\conv Z$, but every $Y\subset\conv Z$ has third coordinate
$1/2$. A small translation $(t,0,-t)$, $t>0$, preserves admissibility but makes
$\conv(Y+h)$ disjoint from $X$. The containment part needs translation restricted
to the parallel space of $\aff Z$, or an interior assumption in the full
mass hyperplane. The direct width proof of strict monotonicity does not use
this containment argument.

\clearpage
\begin{thebibliography}{9}
\bibitem{kline} J. Kline, Properties of the $d$-dimensional earth mover's problem,
\emph{Discrete Applied Mathematics} 265 (2019), 128--141.
\href{https://doi.org/10.1016/j.dam.2019.02.042}{doi:10.1016/j.dam.2019.02.042}.
\bibitem{qst} M. Queyranne, F. Spieksma and F. Tardella, A general class of greedily
solvable linear programs, \emph{Mathematics of Operations Research} 23(4) (1998),
892--908. \href{https://doi.org/10.1287/moor.23.4.892}{doi:10.1287/moor.23.4.892}.
\bibitem{makur} A. Makur and J. Singh, Doeblin coefficients and related measures,
\emph{IEEE Transactions on Information Theory} 70(7) (2024).
\href{https://doi.org/10.1109/TIT.2024.3367856}{doi:10.1109/TIT.2024.3367856}.
Checked manuscript: \href{https://arxiv.org/abs/2309.08475v2}{arXiv:2309.08475v2}.
\bibitem{bln} A. Barvinok, S. J. Lee and I. Novik, Explicit constructions of
centrally symmetric $k$-neighborly polytopes and large strictly antipodal sets,
2012 preprint, \href{https://arxiv.org/abs/1203.6867}{arXiv:1203.6867}.
\bibitem{bt} D. Bryant and P. Tupper, Linear and sublinear diversities,
\href{https://arxiv.org/abs/2412.07092v3}{arXiv:2412.07092v3}, revised March 3, 2026.
\bibitem{mehta} R. Mehta, J. Kline, V. S. Lokhande, G. Fung and V. Singh,
Efficient discrete multi marginal optimal transport regularization,
\emph{International Conference on Learning Representations} (2023).
\href{https://openreview.net/forum?id=R98ZfMt-jE}{OpenReview: R98ZfMt-jE}.
\end{thebibliography}
\end{document}
